\documentclass[aspectratio=169,11pt]{beamer}

% ---- Minimalist theme (matches midway / implementation slides) ----
\usetheme{default}
\usecolortheme{default}
\setbeamertemplate{navigation symbols}{}
\setbeamertemplate{frametitle continuation}{}
\setbeamertemplate{itemize items}[circle]
\setbeamertemplate{enumerate items}[default]

% ---- Font: TeX Gyre Heros ----
\usepackage[T1]{fontenc}
\usepackage{tgheros}
\renewcommand{\familydefault}{\sfdefault}
\usepackage{microtype}

% ---- Packages ----
\usepackage{amsmath}
\usepackage{amssymb}
\usepackage{booktabs}
\usepackage{array}
\usepackage{tikz}
\usepackage{xcolor}
\usepackage{hyperref}
\usetikzlibrary{arrows.meta,positioning,calc,fit}

% ---- Maastricht University colors ----
\definecolor{umdark}{HTML}{001C3D}
\definecolor{umorange}{HTML}{E84E10}
\definecolor{umlight}{HTML}{4A90C4}
\definecolor{umgray}{HTML}{6B7280}
\newcommand{\icite}[1]{{\footnotesize\color{umgray!85}#1}}

% ---- Apply colors to Beamer ----
\setbeamercolor{normal text}{fg=umdark}
\setbeamercolor{frametitle}{fg=umdark}
\setbeamercolor{title}{fg=umdark}
\setbeamercolor{subtitle}{fg=umgray}
\setbeamercolor{author}{fg=umgray}
\setbeamercolor{date}{fg=umgray}
\setbeamercolor{institute}{fg=umorange}
\setbeamercolor{itemize item}{fg=umorange}
\setbeamercolor{itemize subitem}{fg=umlight}
\setbeamercolor{enumerate item}{fg=umorange}
\setbeamercolor{block title}{bg=umdark,fg=white}
\setbeamercolor{block body}{bg=umdark!5,fg=umdark}
\setbeamercolor{block title alerted}{bg=umorange,fg=white}
\setbeamercolor{block body alerted}{bg=umorange!8,fg=umdark}
\setbeamercolor{block title example}{bg=umlight,fg=white}
\setbeamercolor{block body example}{bg=umlight!8,fg=umdark}

% ---- Frametitle style ----
\setbeamertemplate{frametitle}{%
  \vspace{0.35cm}%
  \noindent\hspace*{0pt}%
  \parbox{\textwidth}{%
    {\usebeamerfont{frametitle}\usebeamercolor[fg]{frametitle}\insertframetitle}%
    \vspace{2pt}\\%
    {\color{umorange}\rule{\textwidth}{1.2pt}}%
  }%
  \vspace{-2pt}%
}
\setbeamerfont{frametitle}{size=\large,series=\bfseries}
\setbeamersize{text margin left=0.7cm,text margin right=0.7cm}

% ---- Footline ----
\setbeamertemplate{footline}{%
  \hbox{%
    \begin{beamercolorbox}[wd=\paperwidth,ht=2.2ex,dp=1ex]{footline}%
      \hspace{1em}{\scriptsize\color{umgray}Final Presentation -- Carrier Origin and Steganographic Detectability}%
      \hfill%
      {\scriptsize\color{umgray}\insertframenumber/\inserttotalframenumber}%
      \hspace{1em}%
    \end{beamercolorbox}%
  }%
}

% ---- Metadata ----
\title{Does the Source of Carrier Image Affect\\Steganographic Detectability?}
\subtitle{Final Presentation}
\newcommand{\teamauthors}{Abdul Moiz Akbar \;|\; Malo Coquin \;|\; Daria Gjonbalaj \;|\; Nico Muller-Spath\\Jimena Narvaez del Cid \;|\; David Wicker \;|\; Nikolas Zouros}
\author[Project 2.2 Team]{Project 2.2 Team}
\institute{Department of Advanced Computing Sciences\\Maastricht University}
\date{Project 2.2 \;|\; May 2026}

\begin{document}

% ============================================================
% 1. Opening
% ============================================================
\begin{frame}[plain]
\vfill
\begin{center}
{\color{umorange}\rule{0.42\textwidth}{2pt}}\\[12pt]
{\LARGE\bfseries\color{umdark}\inserttitle}\\[10pt]
{\normalsize\color{umgray}\insertsubtitle}\\[16pt]
{\small\color{umdark}\teamauthors}\\[8pt]
{\small\color{umorange}\insertinstitute}\\[8pt]
{\footnotesize\color{umgray}\insertdate}\\[12pt]
{\color{umorange}\rule{0.42\textwidth}{2pt}}
\end{center}
\vfill
\end{frame}

% ============================================================
% 2. Agenda
% ============================================================
\begin{frame}{Agenda}
\small
\begin{enumerate}
\item Motivation recap
\item Research questions
\item Experimental design and methods
\item Statistical analysis
\item Results across RQ1--RQ5
\item Limitations and contribution
\end{enumerate}
\end{frame}

% ============================================================
% 3. Motivation recap
% ============================================================
\section{Motivation}
\begin{frame}{Motivation Recap}
\small
\begin{itemize}
\item Steganalysis detectability depends on carrier-image statistics, not only on embedding logic \icite{(Petitcolas et al., 1999; Fridrich \& Kodovsky, 2012).}
\item Photographic carriers dominate benchmarks, yet diffusion outputs are now everywhere and carry distinct statistical traces \icite{(Corvi et al., 2023).}
\item Closest prior work studies AI carriers under adaptive embedding and reports $P_E$, not AUC \icite{(M\'ereur et al., 2024).}
\end{itemize}
\vspace{0.1cm}
\begin{alertblock}{Core question}
Do detectors validated on photographs remain equally effective when the carrier is ML-generated, under identical embedding settings?
\end{alertblock}
\end{frame}

% ============================================================
% 4. Research Questions
% ============================================================
\section{Research Questions}
\begin{frame}{Research Questions}
\scriptsize
\begin{columns}[T]
\begin{column}{0.49\textwidth}
\begin{block}{RQ1 (Primary, confirmatory)}
Does carrier source (real vs.\ pooled ML) affect detectability under matched embedding settings?
\end{block}
\begin{block}{RQ2 (Primary, confirmatory)}
Within ML carriers, does the choice of generator (SDXL vs.\ FLUX) change detectability?
\end{block}
\begin{block}{RQ3 (Exploratory)}
Does payload size change the detectability gap between real and ML carriers?
\end{block}
\end{column}
\begin{column}{0.49\textwidth}
\begin{block}{RQ4 (Exploratory)}
Do the spatial (LSB+PNG) and frequency (DCT-LSB+JPEG) branches show different carrier-source effects?
\end{block}
\begin{block}{RQ5 (Verification)}
Does AES-256-CBC payload encryption affect detectability? Expected null result.
\end{block}
\end{column}
\end{columns}
\end{frame}

% ============================================================
% 5. Experimental Design
% ============================================================
\section{Experimental Design}
\begin{frame}{Datasets and Factorial Design}
\small
\begin{columns}[T]
\begin{column}{0.49\textwidth}
\begin{block}{Datasets (1{,}500 images)}
\begin{itemize}
\item Real: 300 COCO + 200 Flickr30k
\item ML-A: SDXL 1.0 (500 caption-matched)
\item ML-B: FLUX.1-schnell (500 caption-matched)
\item Grayscale 512$\times$512, dual storage: PNG and JPEG Q=95
\end{itemize}
\end{block}
\end{column}
\begin{column}{0.49\textwidth}
\begin{block}{$3 \times 2 \times 3 \times 2$ factorial}
\begin{itemize}
\item Source: real / ml\_a / ml\_b
\item Method: LSB / DCT-LSB
\item Payload: low (0.25) / med (0.50) / high (0.75)
\item Encryption: plain / AES-256-CBC
\end{itemize}
\end{block}
\end{column}
\end{columns}
\vspace{0.1cm}
\begin{exampleblock}{Caption-matched, not pixel-matched}
ML generators receive the same captions as real images; subject matter held constant, carrier statistics vary. Per group: 3 covers $\times$ 12 stego conditions $\Rightarrow$ 36 stego images, $\sim$60 detector evaluations.
\end{exampleblock}
\end{frame}

% ============================================================
% 6. Embedding Methods
% ============================================================
\begin{frame}{Embedding Methods}
\small
\begin{columns}[T]
\begin{column}{0.49\textwidth}
\begin{block}{Spatial: LSB substitution}
\begin{itemize}
\item Replace pixel LSB, $k\!=\!1$, row-major
\item Lossless PNG storage required
\item Detected by RS, $\chi^2$, Sample Pairs
\end{itemize}
\end{block}
\end{column}
\begin{column}{0.49\textwidth}
\begin{block}{Frequency: DCT-LSB (JSteg-style)}
\begin{itemize}
\item Flip LSB of \emph{non-zero} quantised AC coefficients
\item Skip DC and $|c|\!=\!1$ (Westfeld-pair stable)
\item JPEG Q=95, single quantisation pass
\end{itemize}
\end{block}
\end{column}
\end{columns}
\vspace{0.1cm}
\begin{exampleblock}{Why classical, not adaptive (HILL / UNIWARD / MiPOD)}
Adaptive schemes adapt to image content, conflating embedding evasion with carrier statistics. Classical fixed-position embedding isolates the carrier-source signal we want to measure.
\end{exampleblock}
\end{frame}

% ============================================================
% 7. Detector Choices
% ============================================================
\begin{frame}{Detector Choices and Mechanisms}
\small
\begin{columns}[T]
\begin{column}{0.49\textwidth}
\begin{block}{Spatial branch}
\begin{itemize}
\item \textbf{RS Analysis} \icite{(Fridrich 2001):} regular/singular pixel-group counts
\item \textbf{$\chi^2$ attack} \icite{(Westfeld 1999):} $\{2k, 2k\!+\!1\}$ pairs-of-values
\item \textbf{Sample Pairs} \icite{(Dumitrescu 2003):} trace multisets on adjacent pairs
\end{itemize}
\end{block}
\end{column}
\begin{column}{0.49\textwidth}
\begin{block}{Frequency branch}
\begin{itemize}
\item \textbf{$\chi^2$ DCT} \icite{(Westfeld 1999):} pairs on quantised AC coefficients
\item \textbf{Calibration $\chi^2$} \icite{(Fridrich 2003):} non-block-aligned crop + recompress reference
\end{itemize}
\end{block}
\end{column}
\end{columns}
\vspace{0.1cm}
\begin{alertblock}{Why training-free}
Classical statistical detectors have no learned parameters. AUC differences between carrier groups reflect carrier statistics, not a training-set mismatch -- critical for an unbiased real-vs-ML comparison.
\end{alertblock}
\end{frame}

% ============================================================
% 7. Pipeline Structure
% ============================================================
\section{Pipeline Structure}
\begin{frame}{Pipeline Structure}
\small
\begin{columns}[T]
\begin{column}{0.49\textwidth}
\begin{block}{Five sequential stages}
\begin{itemize}
\item Data: real downloads + ML generation
\item Manifests: payloads + stego (seeded)
\item Embedding: stego images + PSNR/SSIM/FSIM
\item Detection: 5 classical detectors on cover+stego
\item Aggregation: AUC, contrasts, verdicts, power
\end{itemize}
\end{block}
\end{column}
\begin{column}{0.49\textwidth}
\begin{block}{Run isolation}
\begin{itemize}
\item Each run in its own \texttt{runs/<id>/}
\item Config frozen in \texttt{config.json}
\item Manifests record every condition + seed
\item Re-runs idempotent: done work is skipped
\end{itemize}
\end{block}
\end{column}
\end{columns}
\vspace{0.05cm}
{\scriptsize\color{umgray}\texttt{python run.py <profile>} produces every cover, stego, CSV, figure, and verdict in a self-contained directory; no manual post-processing.}
\end{frame}

% ============================================================
% 8. Web Platform
% ============================================================
\section{Web Platform}
\begin{frame}{Web Platform and Live Monitoring}
\small
\begin{columns}[T]
\begin{column}{0.49\textwidth}
\begin{block}{Architecture}
\begin{itemize}
\item Local Python HTTP server (stdlib only)
\item Single-page browser client, no build step
\item Endpoints: list / launch / preview / kill / detail
\item Launch spawns the runner as a subprocess
\end{itemize}
\end{block}
\end{column}
\begin{column}{0.49\textwidth}
\begin{block}{Live observability}
\begin{itemize}
\item Subprocess stdout via Server-Sent Events
\item Browser terminal panel shows progress live
\item Preview validates config before any compute
\item Run-detail surfaces metrics, figures, gallery
\end{itemize}
\end{block}
\end{column}
\end{columns}
\vspace{0.05cm}
{\scriptsize\color{umgray}Power estimate in the drawer: Hanley--McNeil computed live in JavaScript, colour-coded against the proposal threshold $\delta\!=\!0.05$.}
\end{frame}

% ============================================================
% 9. Statistical Analysis
% ============================================================
\section{Statistical Analysis}
\begin{frame}{Statistical Analysis}
\small
\begin{columns}[T]
\begin{column}{0.49\textwidth}
\begin{block}{Tests}
\begin{itemize}
\item DeLong paired ROC test for correlated AUCs \icite{(DeLong 1988)}
\item Bonferroni--Holm correction across the confirmatory family
\item 95\% CIs for exploratory comparisons (RQ3, RQ4)
\item Equivalence margin $\pm 0.025$ AUC for RQ5
\end{itemize}
\end{block}
\end{column}
\begin{column}{0.49\textwidth}
\begin{block}{Power (Hanley--McNeil)}
\begin{itemize}
\item $\alpha = 0.05$, target power $0.80$
\item Confirmatory family size: 15 strata per RQ
\item Min detectable $\delta_\text{AUC}$ at N=500: RQ1 $\approx 0.057$, RQ2 $\approx 0.066$, RQ5 $\approx 0.050$
\end{itemize}
\end{block}
\end{column}
\end{columns}
\vspace{0.1cm}
\begin{alertblock}{Pragmatic boundary}
N=500 honours the proposal threshold $\delta_\text{min}\!=\!0.05$ for the primary (RQ1) and verification (RQ5) questions. RQ2 is a documented power limitation.
\end{alertblock}
\end{frame}

% ============================================================
% 8. Results: RQ1
% ============================================================
\section{Results}
\begin{frame}{RQ1 -- Real vs.\ Pooled ML}
\small
\begin{block}{Test}
Paired DeLong per stratum (detector $\times$ method $\times$ payload), Holm-adjusted across 15 strata.
\end{block}
\begin{columns}[T]
\begin{column}{0.49\textwidth}
\begin{block}{Headline finding}
\begin{itemize}
\item Pooled $\Delta_\text{AUC}$: [\textit{to fill from} \texttt{exp1\_*\_contrasts.csv}]
\item Significant strata (Holm $p\!<\!0.05$): [\textit{n / 15}]
\item Verdict: [\textit{supported / mixed / not supported}]
\end{itemize}
\end{block}
\end{column}
\begin{column}{0.49\textwidth}
\begin{block}{Interpretation}
\begin{itemize}
\item Direction of the effect, consistent across detectors or not.
\item Where is the gap strongest -- spatial or frequency branch?
\item Cite \texttt{exp1\_real\_vs\_pooled\_ml.png} forest.
\end{itemize}
\end{block}
\end{column}
\end{columns}
\end{frame}

% ============================================================
% 9. Results: RQ2
% ============================================================
\begin{frame}{RQ2 -- SDXL vs.\ FLUX within ML}
\small
\begin{block}{Test}
Paired DeLong, same 15-stratum family, Holm-adjusted. Caption-matched cover pairs across the two generators.
\end{block}
\begin{columns}[T]
\begin{column}{0.49\textwidth}
\begin{block}{Headline finding}
\begin{itemize}
\item Pooled $\Delta_\text{AUC}$: [\textit{from} \texttt{exp2\_*\_contrasts.csv}]
\item Significant strata: [\textit{n / 15}]
\item Verdict: [\textit{supported / mixed / not supported}]
\end{itemize}
\end{block}
\end{column}
\begin{column}{0.49\textwidth}
\begin{block}{Power caveat}
At N=500 the minimum detectable $\delta$ for RQ2 is $\approx 0.066$. Smaller generator-specific effects are below the resolution of this study.
\end{block}
\end{column}
\end{columns}
\end{frame}

% ============================================================
% 10. Results: RQ3 + RQ4
% ============================================================
\begin{frame}{RQ3 and RQ4 -- Exploratory}
\small
\begin{columns}[T]
\begin{column}{0.49\textwidth}
\begin{block}{RQ3: Payload interaction}
\begin{itemize}
\item Real-vs-ML gap by payload level (low / med / high).
\item Trend: [\textit{increasing / non-monotonic / flat}]
\item Cite \texttt{exp3a\_payload\_level\_auc.png} and the source$\times$payload heatmap.
\end{itemize}
\end{block}
\end{column}
\begin{column}{0.49\textwidth}
\begin{block}{RQ4: Branch interaction}
\begin{itemize}
\item Spatial source gap minus frequency source gap, per detector and payload.
\item Pooled $\Delta\Delta$: [\textit{from} \texttt{exp4\_*\_contrasts.csv}]
\item Cite \texttt{exp4\_spatial\_vs\_frequency.png}.
\end{itemize}
\end{block}
\end{column}
\end{columns}
\vspace{0.1cm}
\begin{exampleblock}{Reporting convention}
RQ3 and RQ4 are exploratory: we report point estimates and 95\% CIs, no Holm correction, no significance verdict.
\end{exampleblock}
\end{frame}

% ============================================================
% 11. Results: RQ5
% ============================================================
\begin{frame}{RQ5 -- Encryption Invariance (verification)}
\small
\begin{block}{Test}
Paired DeLong between plain and AES-256-CBC stego on shared cover groups. Equivalence margin $\pm 0.025$ AUC per stratum.
\end{block}
\begin{columns}[T]
\begin{column}{0.49\textwidth}
\begin{block}{Headline finding}
\begin{itemize}
\item Strata within margin: [\textit{n / 45}]
\item Pooled $\Delta_\text{AUC}$: [\textit{near zero, expected}]
\item Verdict: encryption invariant / not supported
\end{itemize}
\end{block}
\end{column}
\begin{column}{0.49\textwidth}
\begin{block}{Why the null matters}
A non-null result would indicate the detectors are reacting to payload structure rather than the embedding distortion itself. The null is the proposal-aligned outcome.
\end{block}
\end{column}
\end{columns}
\end{frame}

% ============================================================
% 12. Limitations
% ============================================================
\section{Limitations}
\begin{frame}{Limitations}
\small
\begin{columns}[T]
\begin{column}{0.49\textwidth}
\begin{block}{Acknowledged deviations}
\begin{itemize}
\item ml\_b: PixArt-$\alpha$ replaced with FLUX.1-schnell (PixArt unavailable on HF Inference API).
\item Power: RQ2 effective $\delta_\text{min} \approx 0.066$ at N=500, not 0.050.
\end{itemize}
\end{block}
\end{column}
\begin{column}{0.49\textwidth}
\begin{block}{Scope-bounded}
\begin{itemize}
\item Two generator families only -- cannot disentangle architecture from training data in RQ2.
\item Grayscale only; no colour information.
\item Two DCT-branch detectors vs.\ three spatial.
\end{itemize}
\end{block}
\end{column}
\end{columns}
\end{frame}

% ============================================================
% 13. Contribution
% ============================================================
\section{Contribution}
\begin{frame}{Contribution and Minimum Deliverable}
\small
\begin{block}{What this study contributes}
\begin{itemize}
\item Caption-matched dataset across three carrier sources under identical embedding.
\item Training-free AUC-based comparison with DeLong and Holm correction.
\item Pre-computed RQ verdicts and pilot-based power analysis as reproducible artefacts.
\end{itemize}
\end{block}
\begin{block}{Minimum deliverable met}
\begin{itemize}
\item One encryption scheme (AES-256-CBC). \checkmark
\item Two embedding approaches: spatial LSB and frequency DCT-LSB. \checkmark
\item Defensible answers to RQ1 and RQ2. \checkmark
\end{itemize}
\end{block}
\end{frame}

% ============================================================
% 14. Closing
% ============================================================
\begin{frame}[plain]
\vfill
\begin{center}
{\color{umorange}\rule{0.32\textwidth}{2pt}}\\[14pt]
{\LARGE\bfseries\color{umdark}Thank you}\\[10pt]
{\large\color{umgray}Questions and Discussion}\\[16pt]
{\small\color{umgray}Code and run artefacts: \texttt{runs/<id>/}}\\[14pt]
{\color{umorange}\rule{0.32\textwidth}{2pt}}
\end{center}
\vfill
\end{frame}

\end{document}
